\subsubsection{長編アニメーション映画『崖の上のポニョ』の構造分析―2編の小さな異郷訪問譚の接合―}
\paragraph*{(1)概要}
スタジオジブリ製作における宮崎駿が関わった長編アニメーション映画には根強い人気がある．その中でも『崖の上のポニョ』を対象として構造分析を行う\cite{ponyo}．物語形式の中でも，「異郷訪問譚」の裏返し構造を用いた分析を行うことで，作品の特性を明らかにする．異郷訪問譚とは日本の代表的な異郷訪問譚である「イザナギの黄泉国訪問譚」や「浦島子」などにもみとめ られるように，古くから使用されてきた物語類型 の一つとも言え，いわゆる成長物語との関連がし ばしば指摘されてきた形式である．
\paragraph*{(2)分析方法}
物語の前半のテーマが後半では逆の順序で出現し，後半では前半の否定ないし対立という形をとる構造を裏返し構造という．その構造を援用する手法により『崖の上のポニョ』を分析していく．
\paragraph*{(3)場面分けの方法}
物語を前半と後半で分けその反対の形をとる場面の名称付けをして当てはめていく．

\subsubsection{物語の訴求分析の理論と手法及び応用事例}
\paragraph*{(1)概要}
物語は，古くから社会や文化の価値観形成に影響を与えてきた重要な要素．神話や昔話だけでなく，現代の娯楽作品（小説，コミック，アニメ，映画など）や広告，ニュース記事にも「物語」が含まれている．これらの物語は時代や文化に応じて変化し，私たちの価値観や意思決定に影響を与える\cite{sokyu}．「物語の訴求力」つまり，物語が持つ魅力や影響力を分析するための理論と手法を構築し，その意義を探ることを目的としている．
\paragraph*{(2)分析方法}
コミックの「進撃の巨人」を対象としてシーケンス分析を行っていく．シーケンス分析とは，物語が持つ「イベントの連鎖」に注目し，その構造を明らかにする分析手法である．この手法は，特に神話や民話などに見られる繰り返しのパターン（スキーマ）を抽出し，それがどのように物語全体の訴求力を形成しているかを検討することを目的とする．シーケンス分析は，物語の中に隠されたテーマや価値観を発見し，それを科学的に解明するための有力な手段となる．
\paragraph*{(3)場面分けの方法}
信頼や疑念といったスキーマを設定し，その内容に当てはまる場面を簡潔な言葉で設定した．

\subsection{結果}


以下で文献調査によって選定された論文の場面の切り方とグルーピング方法についてまとめていく．

\paragraph*{(1)長編アニメーション映画『崖の上のポニョ』の構造分析\cite{ponyo}}

切り方：『崖の上のポニョ』のあらすじを文章単位で区切る

グルーピング：登場人物をマーキングし，裏返しの関係を探る

\paragraph*{(2)神話物語と神話を原型にした現代物語の構造比較\cite{sinwa}}

切り方：物語内の事象（場所/部隊の移動・物語内時間が不連続となる時間経過）

グルーピング：課題の構成要素（発生・目的など）を抽出

\paragraph*{(3)ネットワーク分析を用いたミステリー小説における場面の依存関係の分析\cite{misuteri}}

切り方：物語内の事象（物語内での時間経過・特殊な事象/現象の発生・場面内で行動する人物の増減・主体や主観の変更・物語内での場所の移動）

グルーピング：依存場面を抽出

\paragraph*{(4)ロールプレイングシナリオにおけるクエスト構造に着目した物語分析\cite{rpg}}

切り方：ゲームシナリオ内のクエストで区切る

グルーピング：カテゴリー（発生・経過・結末）を抽出

\paragraph*{(5)物語の訴求構造分析の理論と手法及び応用事例\cite{sokyu}}

切り方：対象作品の話数で区切る

グルーピング：スキーマ（信頼・喜びなど）を抽出

\paragraph*{(6)物語自動生成のための推理小説中の伏線分類\cite{suiri}}

切り方：伏線となる場面を抽出

グルーピング：カテゴリー（エピソード・謎の強調など）を抽出

\paragraph*{(7)その他}

切り方：未記載（鑑賞者の心理体験の分析\cite{anime}・キーワードのみを抽出\cite{onrain}・シナリオソフトの開発\cite{rumiko16}を行い場面分けは主観）

\subsection{視点・出力共通点}
視点では，場面同士の関係を視点としているものが多くみられた．

出力では，各それぞれの分析で設定した要素の頻度を出力していたものが共通点として挙げられる．